% How Universal Are Concept Directions? A Controlled Stress-Test of
% Representation Engineering Under RLHF
% BlackboxNLP 2026 @ EMNLP 2026
%
% (Retitled 2026-07-07: the original working title "The Myth of Universal
% Concept Directions Under RLHF" asserted fragmentation + RLHF amplification,
% which the Qwen results did not support --- refusal is largely universal under
% the prompt-based method, and RLHF slightly *consolidates* directions. The
% question form matches the actual findings and the Repro-track framing.)
%
% Compile with: pdflatex main.tex && bibtex main && pdflatex main.tex && pdflatex main.tex

\documentclass[11pt]{article}

% === ACL 2025 style (download acl.sty separately) ===
% If using the official ACL style file, uncomment the next line:
% \usepackage[review]{acl}
% For now, use a compatible setup:
\usepackage[margin=1in]{geometry}
\usepackage{times}
\usepackage{latexsym}
\usepackage{natbib}

% === Standard packages ===
\usepackage[utf8]{inputenc}
\usepackage[T1]{fontenc}
\usepackage{amsmath}
\usepackage{amssymb}
\usepackage{graphicx}
\usepackage{booktabs}
\usepackage{hyperref}
\usepackage{xcolor}
\usepackage{subcaption}
\usepackage{multirow}

% === Custom macros ===
\newcommand{\cosim}{\text{cos-sim}}
\newcommand{\dmodel}{d_{\text{model}}}
\newcommand{\vdir}{\mathbf{d}}
\newcommand{\vact}{\mathbf{a}}
\newcommand{\R}{\mathbb{R}}

% === Metadata ===
\title{How Universal Are Concept Directions? \\ A Controlled Stress-Test of Representation Engineering Under RLHF}

\author{
  Author A \\
  \textit{Affiliation} \\
  \texttt{email@example.com}
  \and
  Author B \\
  \textit{Affiliation} \\
  \texttt{email@example.com}
}

\date{}

\begin{document}
\maketitle

% ============================================================
\begin{abstract}
% TODO: Write abstract after results are in (Phase 3)
%
% Structure: (1) Motivation --- representation engineering claims universal concept
% directions, (2) Gap --- are these directions truly universal across semantic domains?
% (3) Method --- contrastive prompt pairs, per-domain vs. global direction computation,
% angular dispersion analysis, steering transfer tests, (4) Key findings --- TBD,
% (5) Contribution --- evidence for/against domain fragmentation of concept directions
% under RLHF alignment.
%
\noindent
[Abstract placeholder --- to be written after results are available.]
\end{abstract}

% ============================================================
\section{Introduction}
\label{sec:intro}

% TODO: Person A leads this section
%
% Key points to cover:
% - Representation engineering (Zou et al., 2023) claims single linear directions
%   capture high-level concepts like honesty and refusal in LLM activation space
% - This has practical implications for safety: if concept directions are universal,
%   steering is straightforward; if fragmented, safety interventions may miss domains
% - Our research question: do concept directions generalize uniformly across semantic
%   domains, or fragment into domain-specific sub-directions?
% - RLHF angle: does alignment amplify or reduce fragmentation?
% - Preview of method and findings


% ============================================================
\section{Related Work}
\label{sec:related}

% TODO: Cover the following papers:
%
% Representation Engineering:
% - \citet{zou2023representation} --- canonical concept direction paper
% - \citet{marks2023geometry} --- linear structure of truth in LLM activations
% - \citet{hernandez2023linearity} --- linearity of relation decoding
%
% Activation Steering:
% - \citet{turner2023activation} --- activation addition for steering
% - \citet{li2024inference} --- inference-time intervention
%
% Safety & Alignment:
% - \citet{wei2024assessing} --- brittleness of safety alignment
% - \citet{lieberum2023does} --- scalability of circuit analysis


% ============================================================
\section{Methods}
\label{sec:methods}

Our design asks a simple question of the representation-engineering
framework: if a concept direction is extracted from data spanning several
semantic domains, how well does it agree with directions extracted from each
domain alone? We operationalize this for two concepts --- \textit{honesty}
and \textit{refusal} --- each decomposed into four domains, and measure
agreement both geometrically (angular dispersion) and behaviorally
(activation steering on held-out prompts). Crucially, the design controls
for the confounds that would otherwise masquerade as fragmentation: source
datasets that differ across domains, templated responses, and unbalanced
domain sizes.

\subsection{Contrastive Data}
\label{sec:prompts}

\paragraph{Honesty (primary).}
Domain must be the only factor that varies between the groups we compare, so
our primary experiment draws all of its data from a single source. We
partition the 817 questions of TruthfulQA \citep{lin2022truthfulqa} into
four domains by category --- \textit{factual trivia}, \textit{math},
\textit{politics/opinion}, and \textit{personal advice} --- supplemented
with 30 hand-written arithmetic pairs to populate the math domain. Each item
pairs a question with its best truthful answer (positive) and a plausible
falsehood (negative); activations are read from the concatenation of the
shared question and each response, so the question itself cancels in the
contrast and the measured difference isolates the response.

\paragraph{Refusal.}
We construct the refusal contrast in two independent ways. The
\textbf{prompt-based} condition follows the now-standard recipe of
\citet{arditi2024refusal}: harmful instructions (drawn from AdvBench
\citep{zou2023universal} and MedSafetyBench \citep{han2024medsafetybench},
plus hand-written privacy requests to balance the smallest domain) are
contrasted with harmless instructions, read at the final prompt token, with
no model responses involved. This is our headline refusal condition, as it
is immune to response templating. The \textbf{response-based} condition
pairs the same harmful prompts with compliant versus refusing responses; the
refusals were programmatically diversified so that no phrasing dominates the
negative class. We treat this second condition purely as a robustness check,
since synthetic responses risk encoding surface templates rather than the
target behavior. All pairs carry a framing tag distinguishing
harmful-request refusal from over-refusal of benign requests, and the two
framings are never pooled.

\paragraph{Balance.}
Because a global direction computed over unbalanced domains is dominated by
the largest one, we cap every domain at 120 pairs at extraction time
(prompt-based refusal is balanced at roughly 91--100 prompts per domain by
construction). The full dataset comprises 1{,}720 validated pairs; an
automated validator checks schema conformity, duplicate rates, and response
diversity for every file.

\subsection{Activation Extraction}
\label{sec:extraction}

Our primary models are Qwen~2.5~3B Instruct and its pre-alignment Base
checkpoint \citep{qwen2024qwen25}; Gemma~2~2B \citep{gemmateam2024gemma2}
serves as a pilot. For each prompt we record the post-block residual stream
at every layer (36 for Qwen, 26 for Gemma) at the final token position, in
half precision with NF4 4-bit weight quantization
\citep{dettmers2023qlora} to fit commodity GPUs; we return to the possible
effects of quantization in the Limitations. Extraction metadata (source
files, schemas, pair counts) is written alongside the cached activations so
that every downstream number is traceable to its data.

\subsection{Direction Computation}
\label{sec:directions}

Following \citet{zou2023representation}, the concept direction at layer
$\ell$ is the difference of means
\begin{equation}
  \vdir_\ell \;=\;
  \frac{1}{|P|}\sum_{x \in P} \vact_\ell(x)
  \;-\;
  \frac{1}{|N|}\sum_{x \in N} \vact_\ell(x),
\end{equation}
normalized to unit length, where $P$ and $N$ are the positive and negative
activation sets and $\vact_\ell(x) \in \R^{\dmodel}$ is the residual-stream
activation. Per-domain directions $\vdir_\ell^{(d)}$ use only pairs from
domain $d$; the global direction $\vdir_\ell^{(g)}$ pools all domains after
balancing.

\subsection{Angular Dispersion and Statistics}
\label{sec:angular}

Fragmentation at layer $\ell$ is summarized by the mean cosine similarity
between each per-domain direction and the global direction,
$\overline{\cosim}_\ell = \tfrac{1}{4}\sum_d
\cos\!\big(\vdir_\ell^{(d)}, \vdir_\ell^{(g)}\big)$,
together with the corresponding mean angle. A value near~1 indicates a
single shared axis; lower values indicate domain-specific structure. We also
report the full pairwise domain-by-domain cosine matrix at the most
fragmented layer, excluding the two embedding-adjacent layers, whose
directions reflect token statistics rather than concept representations.
Confidence intervals are obtained by a \emph{pair-level} bootstrap: in each
of 1{,}000 resamples we redraw prompt pairs with replacement within every
domain, recompute all per-domain and global directions from scratch, and
recompute the dispersion statistic. This propagates sampling variability in
the underlying data, rather than resampling the four domain cosines
themselves. Because early layers of a causal language model are dominated by
lexical features, we additionally report the mean cosine over the final
third of the network, where representations feed the unembedding and
behavior is read out.

\subsection{Behavioral Steering}
\label{sec:steering-method}

Geometry alone cannot establish that two directions are functionally
interchangeable, so we test transfer behaviorally with activation addition
\citep{turner2023activation,li2024inference}. For each refusal domain we
hold out 20 harmful prompts that were never used for direction computation,
format them with the model's chat template, and generate 64 tokens greedily
under three conditions: unsteered, steered with the global direction, and
steered with the domain's own direction (plus each other domain's direction
in a cross-domain sweep). The direction, scaled by a coefficient of 4.0, is
added to the residual stream at layer 18 at every token position; its sign
points from refusal toward compliance, so effective steering should
\emph{reduce} the refusal rate on harmful prompts. Refusals are detected by
keyword matching over a standard set of refusal markers. If concept
directions fragment, a domain's own direction should outperform the global
one on that domain's prompts; if they are universal, the two should be
indistinguishable.

\subsection{RLHF Comparison}
\label{sec:rlhf-method}

To isolate the effect of alignment we repeat the entire pipeline ---
extraction, direction computation, and dispersion analysis --- on the Base
checkpoint, which shares the pre-training of the Instruct model but has
undergone no preference tuning, and compare
$\overline{\cosim}_\ell$ layer by layer across the two checkpoints. The
pre-registered hypothesis was that alignment \emph{increases} dispersion. As
a sanity check that the checkpoints genuinely differ, we verify that
per-domain cosine differences between Base and Instruct grow with depth, the
expected signature of a fine-tune.


% ============================================================
\section{Experimental Setup}
\label{sec:setup}

\begin{table}[t]
  \centering
  \small
  \begin{tabular}{lll}
    \toprule
    & \textbf{Honesty} & \textbf{Refusal} \\
    \midrule
    Source & TruthfulQA & AdvBench, MedSafety- \\
           &            & Bench, hand-written \\
    Domains & 4 & 4 \\
    Pairs/domain & $\le$120 (math 59) & 91--120 \\
    Contrast & true vs.\ false & harmful vs.\ harmless \\
             & answer & prompt (+ resp.\ variant) \\
    \bottomrule
  \end{tabular}
  \caption{Data summary. Both concepts are additionally extracted from the
  Qwen~2.5~3B Base checkpoint for the RLHF comparison.}
  \label{tab:setup}
\end{table}

All experiments ran on single T4/P100 GPUs (roughly 25 GPU-hours in total);
analysis is CPU-only and fully scripted, from cached activations to the
figures in this paper. We use greedy decoding throughout the steering
experiments and fix all random seeds; the mock-data mode of the released
pipeline reproduces every table in this paper end-to-end on synthetic
activations before any GPU is touched. Code, data manifests, and validation
tooling are released with the paper.


% ============================================================
\section{Results}
\label{sec:results}

% NOTE: copy figure PDFs into paper/figures/ before compiling, e.g.
%   cp new_results_final/figures/qwen/*.pdf paper/figures/
%   cp "new_results_final/figures/qwen respbased/"*.pdf paper/figures/respbased_...

\begin{table}[t]
  \centering
  \small
  \begin{tabular}{lccc}
    \toprule
    \textbf{Condition} & \textbf{Min $\overline{\cosim}$} & \textbf{Layer} & \textbf{Late third} \\
    \midrule
    Refusal (prompt)    & 0.837 & 0  & 0.928 \\
    Refusal (response)  & 0.854 & 0  & 0.904 \\
    Honesty             & 0.522 & 5  & 0.694 \\
    \midrule
    \multicolumn{4}{l}{\textit{Base checkpoint}} \\
    Refusal (prompt)    & 0.834 & 0  & 0.911 \\
    Honesty             & 0.522 & 5  & 0.667 \\
    \bottomrule
  \end{tabular}
  \caption{Angular dispersion on Qwen~2.5~3B: minimum mean cosine between
  per-domain and global directions across layers, and the mean over the final
  third of the network. Refusal directions are close to a single axis
  throughout; honesty directions diverge sharply in early-mid layers and only
  partially reconverge.}
  \label{tab:dispersion}
\end{table}

\subsection{Refusal Directions Are Largely Universal}
\label{sec:results-refusal}

Contrary to our first pre-registered hypothesis, refusal shows little domain
structure. Under the prompt-based condition, the mean cosine between
per-domain and global directions never drops below 0.837 at any of the 36
layers of Qwen~2.5~3B Instruct, and averages 0.928 over the final third of
the network (Table~\ref{tab:dispersion}). The response-based robustness
condition agrees: despite an entirely different contrast construction, mean
cosines span 0.854--0.927 (pair-level 95\% CI at the most dispersed layer:
$[0.791, 0.885]$), with a late-layer mean of 0.904. Whatever the network
uses to decide refusal, it points in essentially the same direction whether
the request concerns violence, illegal activity, medical ethics, or privacy
--- consistent with the single-direction account of
\citet{arditi2024refusal}, and now shown to be stable across two
operationalizations of the concept.

\subsection{Honesty Fragments in Early and Middle Layers}
\label{sec:results-honesty}

Honesty behaves differently. The mean cosine to the global direction falls
to 0.522 at layer~5 --- a mean angle of $57^\circ$ --- and the domain
structure at that depth is stark: the math direction is essentially
orthogonal to every other domain (pairwise cosines between $-0.10$ and
$-0.03$), and even the best-aligned pair, factual trivia and personal
advice, reaches only 0.34. Dispersion then decays with depth, but only
partially: the final-third mean of 0.694 remains well below the
corresponding refusal values. Universality is thus not a property of concept
directions per se, but of particular concepts at particular depths. We note
that early-layer divergence may partly reflect lexical differences between
domains (arithmetic text simply looks different from trivia questions), which
is one reason we insist on the behavioral test below; the late-layer gap
between honesty and refusal, however, survives this caveat.

\subsection{Steering Confirms the Geometry}
\label{sec:results-steering}

\begin{table}[t]
  \centering
  \small
  \begin{tabular}{lccc}
    \toprule
    \textbf{Domain} & \textbf{Baseline} & \textbf{Global $\Delta$} & \textbf{Own $\Delta$} \\
    \midrule
    violence          & 0.95 & \phantom{$-$}0.00 & \phantom{$-$}0.00 \\
    illegal activity  & 0.90 & $-$0.05 & $-$0.05 \\
    medical/legal     & 0.50 & $-$0.15 & \phantom{$-$}0.00 \\
    privacy           & 0.50 & $-$0.40 & $-$0.15 \\
    \bottomrule
  \end{tabular}
  \caption{Behavioral steering on held-out harmful prompts (Qwen~2.5~3B
  Instruct, layer 18, coefficient 4.0, $n{=}20$ per domain). Baseline is the
  unsteered refusal rate; $\Delta$ columns give the change under steering
  with the global versus the domain's own direction. The global direction is
  never worse than the domain-specific one.}
  \label{tab:steering}
\end{table}

The second hypothesis --- that global steering underperforms domain-specific
steering --- also fails, in the direction favorable to universality
(Table~\ref{tab:steering}). In every domain the global direction shifts
behavior at least as much as the domain's own direction, and in two domains
strictly more (medical/legal: $-0.15$ vs.\ $0.00$; privacy: $-0.40$ vs.\
$-0.15$). Cross-domain directions transfer as well: privacy prompts are
steered as strongly by the violence and illegal-activity directions
($-0.40$) as by the global one. We attribute the weaker own-domain effects
to estimation noise --- each per-domain direction is fit on roughly a
quarter of the data --- which is itself an argument for pooling. Two
patterns deserve note: absolute effects are modest at this coefficient, and
the violence domain is behaviorally immovable (baseline refusal 0.95,
unchanged under every direction), suggesting refusal of violent requests is
enforced more redundantly than a single rank-one intervention can undo.

\subsection{RLHF Does Not Amplify Fragmentation}
\label{sec:results-rlhf}

Our third hypothesis predicted that alignment would increase dispersion. The
opposite holds, weakly: the Instruct checkpoint is slightly \emph{more}
universal than Base, with a mean per-layer gain in
$\overline{\cosim}_\ell$ of $+0.011$ for honesty and $+0.007$ for refusal,
and the gap widening toward the output (honesty at layer 35: 0.690 Base
vs.\ 0.769 Instruct). Only 10 of 36 layers move in the predicted direction,
none substantially. Base--Instruct differences in individual per-domain
cosines grow with depth (up to 0.20), the expected signature of a fine-tune,
confirming that the two checkpoints genuinely diverge --- alignment simply
does not diverge them along the axis we hypothesized. If anything,
preference tuning appears to \emph{consolidate} concept representations.

\subsection{A Cautionary Non-Replication}
\label{sec:results-pilot}

Our pilot on Gemma~2~2B, run on the response-based refusal data, showed what
looked like strong fragmentation: a minimum mean cosine of 0.559 and
near-orthogonal domain pairs. None of this survives on Qwen~2.5~3B --- the
same data recipe yields cosines of 0.85 and above (\S\ref{sec:results-refusal}).
We initially suspected dataset artifacts (mixed sources, templated
responses), but the prompt-based and response-based conditions agree on
Qwen, ruling that explanation out. The fragmentation was a property of the
smaller model, not of the data. We report this as a caution for the
literature: single-model fragmentation findings, including ours on Gemma,
should not be read as facts about concept directions in general.


% ============================================================
\section{Discussion}
\label{sec:discussion}

% TODO: Interpret results
%
% Key questions to address:
% - Do we find fragmentation? If so, in which layers/domains?
% - What does this mean for representation engineering as a safety tool?
% - Does RLHF help or hurt universality?
% - Implications for activation steering approaches
% - Connection to prior work


% ============================================================
\section{Limitations}
\label{sec:limitations}

% TODO: Cover:
% - Model scale (3B parameters --- may not generalize to larger models)
% - Only 2 concepts (refusal, honesty)
% - Prompt pair quality and coverage
% - Last-token extraction may miss information
% - Linear direction assumption


% ============================================================
\section{Conclusion}
\label{sec:conclusion}

% TODO: Write after results are in


% ============================================================
\section*{Acknowledgments}

% TODO: Add acknowledgments


% ============================================================
\bibliographystyle{acl_natbib}
\bibliography{references}

% ============================================================
% Appendix (optional)
\appendix

\section{Prompt Pair Examples}
\label{app:prompts}

% TODO: Include representative examples from each domain

\section{Full Layer-wise Results}
\label{app:layers}

% TODO: Include detailed per-layer results tables

\end{document}
